
\sectiontoc{1. Introduction}

The success of Maxwell's equations has led to electrodynamics 
being normally formulated in terms of fields that have degrees 
of freedom independent of the particles in them. However, Gauss 
suggested that an action--at--a--distance theory in which the 
action travelled at a finite velocity might be possible. 
This idea was developed by Wheeler and Feynman~\cite{2,3} who derived 
their theory from an action principle that involved only direct 
interactions between pairs of particles. A feature of this theory 
was that the `pseudo'--fields introduced are the half--retarded 
plus half--advanced fields calculated from the world--lines of 
the particles. 

However, Wheeler and Feynman, and, in a different way, Hogarth~(3), 
were able to show that, provided certain cosmological conditions 
were satisfied, these fields could combine to give the observed field. 
Hoyle and Narlikar~\cite{4} extended the theory to general space--times 
and obtained similar theories for their `C'--field~\cite{5} and for the 
gravitational field~~\cite{6}. It is with these theories that this chapter 
is concerned.

It will be shown that in an expanding universe the advanced 
fields are infinite, and the retarded fields finite. 
This is because, unlike electric charges, all masses have the same sign.

\sectiontoc{The Boundary Condition}

Hoyle and Narlikar derive their theory from the action:
\[
A = \sum_{a \ne b} \sum \iint G(a,b)\, da\, db ,
\]
where the integration is over the world-lines of particles a, b...
In this expression G is a Green function
that satisfies the wave equation: 

\[
G(x,x')_{;ij}\, g^{ij} + \tfrac{1}{6} R\, G(x,x')
= \frac{\delta^{4}(x,x')}{\sqrt{-g}} \, .
\]


where \(S\) is the determinant of \(g_{ij}\).
Since the double sum in the action A is symmetrical between all pailrs of
particles a, b, only that part of G(a,b) that is symmetrical between a and b will contribute to the action
i.e. the action can be written

\[
A = \sum_{a \ne b} \sum \iint G^{*}(a,b)\, da\, db
\]

\noindent
where \quad 
\( G^{*}(a,b) = \tfrac{1}{2} G(a,b) + \tfrac{1}{2} G(b,a) \).

\medskip

Thus \quad \( G^{*} \) \quad must be the time-symmetric Green function, and can be written:

\[
G^{*} = \tfrac{1}{2} G_{\text{ret}} + \tfrac{1}{2} G_{\text{adv}}
\qquad \text{where } G_{\text{ret}}
\]
and \(G_{\text{adv}}\) are the retarded and advanced Green functions.
By requiring that the action be stationary under variations
of \(g_{ij}\), Hoyle and Narlikar obtain the field-equations:

\[
\Biggl[
\sum_{a \ne b} \sum \tfrac{1}{6} m^{(a)}(x)\, m^{(b)}(x)
\Biggr]
\left( R_{ik} - \tfrac{1}{2} g_{ik} R \right)
\]

\[
= - T_{ik}
+ \sum_{a \ne b} \sum \tfrac{1}{3} m^{(a)}
\Bigl[
g_{ik} m^{(b)r}{}_{;r}
- m^{(b)}_{;ik}
\Bigr]
+ 2 \bigl( m^{(a)}_{il}\, m^{(b)}_{j}{}^{l} \bigr)
- \tfrac{1}{4} g_{ik}\, m^{(a);r}\, m^{(b)}_{j;r}
\]
where \quad 
\( m^{(a)}(x) = \int G^{*}(x,a)\, da \).

\medskip

However, as a consequence of the particular choice of Green function, 
the contraction of the field--equations is satisfied identically. 
There are thus only 9 equations for the 10 components of \( g_{ij} \) 
and the system is indeterminate.

\medskip

Hoyle and Narlikar therefore impose 
\(\sum m^{(a)} = m_{0} = \text{const.}\), as the tenth equation. 
By then making the `smooth--fluid' approximation, that is by putting

\[
\sum_{a \ne b} \sum m^{(a)} m^{(b)} \approx m_{0}^{2},
\]

they obtain the Einstein field--equations:

\[
\tfrac{1}{6} m_{0}^{2}\left( R_{ik} - \tfrac{1}{2} R g_{ik} \right) 
= -T_{ik}.
\]

\medskip

There is an important difference, however, between these field--equations 
in the direct--particle interaction theory and in the usual general theory 
of relativity. In the general theory of relativity, any metric that satisfies the the field--equations is admissible, but in the direct--particle 
interaction theory only those solutions of the field--equations 
are admissible that satisfy the additional requirement:

\[
m_{0}(x) = \sum m^{(a)}(x) = \sum \int G^{*}(x,a)\, da
= \tfrac{1}{2} \sum \int G_{\text{ret}}(x,a)\, da 
+ \tfrac{1}{2} \sum \int G_{\text{adv}}(x,a)\, da
\]

\medskip

This requirement is highly restrictive; it will be shown that it 
is not satisfied for the cosmological solutions of the Einstein 
field--equations, and it appears that it cannot be satisfied for 
any models of the universe that either contain an infinite amount 
of matter or undergo infinite expansion.

\medskip

The difficulty is similar to that occurring in Newtonian theory 
when it is recognized that the universe might be infinite.  

The Newtonian potential \(\phi\) obeys the equation:

\[
\square \phi = - \kappa \rho \qquad (\rho > 0),
\]

where \(\rho\) is the density.
In an infinite static universe, \(\phi\) would be infinite, since 
the source always has the same sign. The difficulty was resolved 
when it was realized that the universe was expanding, since 
in an expanding universe the retarded solution of the above 
equation is finite by a sort of `red--shift' effect. The 
advanced solution will be infinite by a `blue--shift' effect. 
This is unimportant in Newtonian theory, since one is free 
to choose the solution of the equation and so may ignore the 
infinite advanced solution and take simply the finite 
retarded solution.

\medskip

Similarly in the direct--particle interaction theory the 
\(m\) field satisfies the equation:

\[
\square m - \tfrac{1}{6} R m = \mathcal{N} \qquad (\mathcal{N} > 0),
\]

where \(\mathcal{N}\) is the density of world--lines of particles. 
As in the Newtonian case, one may expect that the effect of the 
expansion of the universe will be to make the retarded solution 
finite and the advanced solution infinite. However, one is 
now not free to choose the finite retarded solution, for the 
equation is derived from a direct--particle interaction action--principle 
symmetric between pairs of particles, and one must choose for \(m\) 
half the sum of the retarded and advanced solutions. We would expect 
this to be infinite, and this is shown to be so in the next section.
\sectiontoc{The Cosmological Solutions}

The Robertson--Walker cosmological metrics have the form
\[
ds^{2} = dt^{2} - R^{2}(t) \left[ \frac{dr^{2}}{1 - \kappa r^{2}} 
+ r^{2}\bigl(d\theta^{2} + \sin^{2}\theta \, d\phi^{2}\bigr) \right].
\]

Since they are conformally flat, one can choose coordinates in which they become
\[
ds^{2} = \Omega^{2} \bigl[ d\tau^{2} - d\rho^{2} - \rho^{2}(d\theta^{2} + \sin^{2}\theta \, d\phi^{2}) \bigr]
= \Omega^{2} \eta_{ab}\, dx^{a} dx^{b},
\]
where \(\eta_{ab}\) is the flat--space metric tensor and
\[
\Omega = \Omega(\tau,\rho) = \frac{R(t)}{\sqrt{[t + \tfrac{1}{4}\kappa(\tau+\rho)^{2}][t + \tfrac{1}{4}\kappa(\tau-\rho)^{2}]}}. \tag{7}
\]

For example, for the Einstein--de Sitter universe
\[
\kappa = 0, \quad 
R(t) = \left(\tfrac{t}{T}\right)^{2/3}, 
\qquad (0 < t < \infty),
\]
\[
\Omega = R \cdot \left(\tfrac{\tau}{T}\right)^{2/3}, 
\qquad (0 < \tau < \infty),
\]
\[
\rho = \tau^{2/3}\, \xi^{1/3}.
\]

For the steady--state (de Sitter) universe
\[
\kappa = 0, \quad R(t) = e^{t/T}, 
\qquad (-\infty < t < \infty).
\]

\[
\Omega = R \cdot \frac{1}{\tau}, 
\qquad (-\infty < t < \infty),
\]
\[
r = \rho \left( \frac{\tau}{T} - T e^{t/T} \right).
\]

The Green function \( G^{*}(a,b) \) obeys the equation
\[
\square G^{*}(a,b) - \tfrac{1}{6} R\, G^{*}(a,b) 
= \frac{\delta^{4}(a,b)}{\sqrt{-g}}.
\]

From this it follows that
\[
\frac{1}{\Omega^{2}} \frac{\partial}{\partial x^{a}}
\left( \Omega^{2} \eta^{ab} \frac{\partial}{\partial x^{b}} G^{*} \right)
= \frac{1}{\Omega^{2}} \frac{\partial}{\partial x^{a}}
\left( \eta^{ab} \frac{\partial}{\partial x^{b}} S \right)
= \Omega^{-3}\, \delta^{4}(a,b).
\]

If we let \( G^{*} = \Omega^{-1} S \), then
\[
\Omega^{2} \frac{\partial}{\partial x^{a}}
\left( \eta^{ab} \frac{\partial}{\partial x^{b}} S \right)
= \delta^{4}(a,b).
\]

This is simply the flat--space Green function equation, and hence
\[
G^{*}(t_{1},0; t_{2},\rho) 
= \frac{\Omega(t_{1})}{8\pi \, \Omega(t_{2}) \, \rho}
\Bigl[ \delta(\rho - (t_{2}-t_{1})) + \delta(\rho + (t_{2}-t_{1})) \Bigr].
\]

The \(m\)--field is given by
\[
m(t,x) = \int G^{*} \mathcal{N}\, \sqrt{-g}\, dx
= \tfrac{1}{2}(m_{\text{ret}} + m_{\text{adv}}).
\]

For universes without creation (e.g.\ the Einstein--de Sitter universe),
\[
\mathcal{N} = R^{3} n, \qquad n = \text{const.}
\]

For Universes with creation (steady state) 
\(\mathcal{N} = n, \quad n = \text{const.}\)

\[
M_{\text{adv}}(\tau_{1}) = \Omega^{-1}(\tau_{1}) 
\int \frac{\mathcal{N}\,\Omega^{3}(\tau_{2})}{4\pi r}\, 4\pi r^{2}\, d\tau_{2},
\]

where the integration is over the future light cone. 
This will normally be infinite in an expanding universe, e.g.\ in the Einstein--de Sitter universe:

\[
M_{\text{adv}}(\tau_{1}) 
= \left(\frac{\tau_{1}}{T}\right)^{-2} 
\int_{\tau_{1}}^{\infty} n \left(\tfrac{\tau_{2}}{T}\right)^{2}\, d\tau_{2}
= \infty.
\]

In the steady--state universe
\[
M_{\text{adv}}(\tau_{1}) 
= \left(\tfrac{-T}{\tau_{1}}\right)^{-1} 
\int_{\tau_{1}}^{0} n \left(\tfrac{-T}{\tau_{2}}\right)^{3} (\tau_{2}-\tau_{1})\, d\tau_{2}
= \infty.
\]

\medskip

By contrast, on the other hand, we have
\[
M_{\text{ret}}(\tau_{1}) = \Omega^{-1}(\tau_{1}) 
\int \frac{\mathcal{N}\,\Omega^{3}}{4\pi r}\, 4\pi r^{2}\, dr,
\]

where the integration is over the past light cone. This will Normally be finite, e.g.\ in the Einstein--de Sitter universe
\[
M_{\text{ret}}(\tau_{1}) 
= \left(\frac{\tau_{1}}{T}\right)^{-2} 
\int_{0}^{\tau_{1}} n \left(\tfrac{\tau_{2}}{T}\right)^{2} (\tau_{2}-\tau_{1})\, d\tau_{2}
= \tfrac{1}{2} n T^{2}.
\]

While in the steady--state universe
\[
M_{\text{ret}}(\tau_{1}) 
= \left(\tfrac{-T}{\tau_{1}}\right)^{-1} 
\int_{-\infty}^{\tau_{1}} n \left(\tfrac{-T}{\tau_{2}}\right)^{3} (\tau_{2}-\tau_{1})\, d\tau_{2}
= \tfrac{1}{2} n T^{2}.
\]

Thus it can be seen that the solution \(m = \text{const.}\) of the equation
\[
\square m + \tfrac{1}{6} R m = \mathcal{N}
\]
is not, in a cosmological metric, the half--advanced plus half--retarded solution 
since this would be infinite. In fact, in the case of the Einstein--de Sitter 
and steady--state metrics, it is the pure retarded solution.

\sectiontoc{The `C'--Field}

Hoyle and Narlikar derive their direct--particle interaction theory of the 
`C'--field from the action
\[
A = \sum_{a \ne b} \iint \hat{G}(a,b)\, i_{a} \, k_{b}\, da^{4} db^{4}.
\]

where the suffixes \(a, b\) refer to differentiation of 
\(\hat{G}(a,b)\) on the world--lines of \(a, b\) respectively.  
\(\hat{G}\) is a Green function obeying the equation

\[
\square \hat{G}(x,x') = \frac{\delta^{4}(x,x')}{\sqrt{-g}}.
\]

We define the `C'--field by
\[
C(x) = \sum \hat{G}(x,a)\, i_{a}\, da^{i},
\]

and the matter--current \(J^{\kappa}\) by
\[
J^{\kappa}(y) = \sum \int \delta^{4}(y,b)\, db^{\kappa}.
\]

Then
\[
C(x) = \int \hat{G}(x,y)\, J^{\kappa}(y)_{;\kappa}\, \sqrt{-g}\, dx',
\]

so that
\[
\square C = J^{\kappa}{}_{;\kappa}.
\]

\medskip

We thus see that the sources of the `C'--field are the places 
where matter is created or destroyed.  

As in the case of the \(m\)--field, the Green function must be 
time--symmetric, that is
\[
\hat{G}(a,b) = \tfrac{1}{2}\, \hat{G}_{\text{ret}}(a,b) 
+ \tfrac{1}{2}\, \hat{G}_{\text{adv}}(a,b)
\]
Hoyle and Narlikar claim that if the action of the 
`C'--field is included along with the action of the `m'--field, 
a universe will be obtained that approximates to the steady--state 
universe on a large scale although there may be local irregularities. 
In this universe, the value of \(C\) will be finite and its gradient 
time--like and of unit magnitude.

\medskip

Given this universe, we may check it for consistency by calculating 
the advanced and retarded `C'--fields and finding if their sum is finite. 
We shall not do this directly but will show that the advanced field 
is infinite while the retarded field is finite.

\medskip

Consider a region in space--time bounded by a three--dimensional 
space--like hypersurface \(\mathcal{D}\) at the present time, 
and the past light cone \(\Sigma\) of some point \(P\) to the future of \(\mathcal{D}\).

By Gauss's theorem
\[
\int_{V} \square C \, \sqrt{-g}\, dx^{4} 
= \int_{\Sigma + \mathcal{D}} \frac{\partial C}{\partial n}\, dS
= \int J^{\kappa}{}_{;\kappa}\, \sqrt{-g}\, dx^{4}.
\]

Let the advanced field produced by sources within \(V\) be \(C'\).  
Then \(C'\) and \(\frac{\partial C'}{\partial n}\) will be zero on \(\Sigma\), and hence


\[
\int_{V} J^{\kappa}{}_{;\kappa} \, \sqrt{-g}\, dx^{4} 
= \int_{\mathcal{D}} \frac{\partial C'}{\partial n}\, dS.
\]

But \(J^{\kappa}{}_{;\kappa}\) is the rate of creation of matter \(= n\) (const.) 
in the steady--state universe, and hence
\[
\int_{\mathcal{D}} \frac{\partial C'}{\partial n}\, dS = nV.
\]

As the point \(P\) is taken further into the future, the volume 
of the region \(V\) tends to infinity. However, the area of the 
hypersurface \(\mathcal{D}\) tends to a finite limit owing to horizon 
effects. Therefore the gradient \(\tfrac{\partial C'}{\partial n}\) must be infinite.  
A similar calculation shows the gradient of the retarded 
field to be finite. Their sums cannot therefore give the 
field of unit gradient required by the Hoyle--Narlikar theory.

\medskip

It is worth noting that this result was obtained without 
assumptions of a smooth distribution of matter or of conformal flatness.
\sectiontoc{Conclusion}

It is one of the weaknesses of the Einstein theory of 
relativity that although it furnishes field equations it does 
not provide boundary conditions for them. Thus it does not 
give a unique model for the universe but allows a whole series 
of models. Clearly a theory that provided boundary conditions 
and thus restricted the possible solutions would be very 
attractive. The Hoyle--Narlikar theory does just that (the 
requirement that \( m = \tfrac{1}{2} m_{\text{ret}} + \tfrac{1}{2} m_{\text{adv}} \) 
is equivalent to a boundary condition). Unfortunately, as we 
have seen above, this condition excludes those models that 
seem to correspond to the actual universe, namely the 
Robertson--Walker models.

\medskip

The calculations given above have considered the universe 
as being filled with a uniform distribution of matter. This 
is legitimate if we are able to make the `smooth--fluid' 
approximation to obtain the Einstein equations. Alternatively, 
if this approximation is invalid, it cannot be said that the 
theory yields the Einstein equations.

\medskip

It might possibly be that local irregularities could make 
\(m_{\text{adv}}\) finite, but this has certainly not been demonstrated 
and seems unlikely in view of the fact that, in the Hoyle--Narlikar 
direct--particle interaction theory of their `C'--field,

which is derived from a very similar action--principle, it can 
be shown without assuming a smooth distribution that the 
advanced `C' field will be infinite in an expanding universe 
with creation.

\medskip

The reason that it is possible to formulate a direct--particle 
interaction theory of electrodynamics that does not encounter 
this difficulty of having the advanced solution infinite is 
that in electrodynamics there are equal numbers of sources 
of positive and negative sign. Their fields can cancel each 
other out and the total field can be zero apart from local 
irregularities. This suggests that a possible way to save the 
Hoyle--Narlikar theory would be to allow masses of both positive 
and negative sign. The action would be
\[
A = \sum_{a \ne b} q_{a} q_{b} \iint G^{*}(a,b)\, da\, db, 
\qquad (q_{a}, q_{b} = \pm 1),
\]
where \(q_{a}, q_{b}\) are gravitational charges analogous to 
electric charges. Particles of positive \(q\) in a positive 
`m'--field and particles of negative \(q\) in a negative 
`m'--field would have the normal gravitational properties, 
that is, they would have positive gravitational and inertial masses.

A particle of negative \(q\) in a positive `m'--field would 
still follow a geodesic. Therefore it would be attracted by 
a particle of positive \(q\). Its own gravitational effect 
however would be to repel all other particles. Thus it would 
have the properties of the negative mass described by Bondi~(8), 
that is, negative gravitational mass and negative inertial mass.

\medskip

Since there does not seem to be any matter having these 
properties in our region of space (where \(m = \text{const.} > 0\)) 
there must clearly be separation on a very large scale. It 
would not be possible to identify particles of negative \(q\) 
with antimatter, since it is known that antimatter has positive 
inertial mass. However, the introduction of negative masses 
would probably raise more difficulties than it would solve.






\sectiontoc{References}

\begin{enumerate}
\item J. A. Wheeler and R. P. Feynman, 
\textit{Rev. Mod. Phys.} \textbf{17}, 157 (1945).

\item J. A. Wheeler and R. P. Feynman, 
\textit{Rev. Mod. Phys.} \textbf{21}, 425 (1949).

\item J. E. Hogarth, 
\textit{Proc. Roy. Soc. A} \textbf{267}, 365 (1962).

\item F. Hoyle and J. V. Narlikar, 
\textit{Proc. Roy. Soc. A} \textbf{277}, 1 (1964a).

\item F. Hoyle and J. V. Narlikar, 
\textit{Proc. Roy. Soc. A} \textbf{282}, 178 (1964b).

\item F. Hoyle and J. V. Narlikar, 
\textit{Proc. Roy. Soc. A} \textbf{282}, 191 (1964c).

\item L. Infeld and A. Schild, 
\textit{Phys. Rev.} \textbf{68}, 250 (1945).

\item H. Bondi, 
\textit{Rev. Mod. Phys.} \textbf{29}, 423 (1957).
\end{enumerate}

